\section{Introduction}
\label{sec:introduction}

Kernel ridge regression (KRR) combines a flexible function class with a
quadratic objective, but its usual data-space formulation contains a dense
\(N\times N\) kernel matrix. A direct solve costs \(O(N^3)\) operations and
\(O(N^2)\) memory. Conjugate gradient (CG) removes the factorization, yet a
dense kernel matrix-vector product still costs \(O(N^2)\)
\cite{aronszajn1950theory,kimeldorf1971some,golub2013matrix,trefethen1997numerical}.
Large-scale KRR therefore needs both a cheaper representation of the kernel
matrix and an iterative system that converges in relatively few steps.

For low-dimensional translation-invariant kernels, an equispaced Fourier
expansion provides the first part. It replaces the data-space system by a
weight-space system with \(M\) Fourier coefficients, where \(M\ll N\) in the
problems considered here. The coefficient matrix has the form
\(A=DGD+\lambda I_M\): \(D\) contains spectral weights and \(G\) is multilevel
Toeplitz. The data enter only through one-time type-1 nonuniform fast Fourier
transform (NUFFT) computations of the Toeplitz generator and right-hand side.
After \(O(N+M\log M)\) setup, a product with \(A\) costs
\(O(M\log M)\), with no direct dependence on \(N\) for fixed \(M\),
\cite{greengard2025efgp,barnett2024uniform}. This separation is especially
useful when the number of observations is much larger than the Fourier grid.

The Fourier approximation and the iterative solution address different
errors. The grid spacing and width determine the finite feature space and the
approximation to the kernel, whereas the residual tolerance determines how
accurately the resulting linear system is solved. Enlarging the grid may
improve the kernel approximation, but it also increases \(M\) and the cost of
each product. A preconditioner instead leaves the exact solution of the fixed
Fourier-discretized system unchanged and changes the work needed to reach the
residual tolerance.  Our algorithmic contribution addresses this second cost.
The empirical comparison is broader and keeps the model-level and solver-level
questions separate: end-to-end KRR pipelines are compared with their own setup
costs, whereas the mechanism study fixes the Fourier model and varies only the
solver.

The low-dimensional restriction is substantive. For a tensor grid with \(2m+1\)
indices per coordinate, \(M=(2m+1)^d\), so neither the Toeplitz product nor the
proposed box removes exponential growth with dimension \(d\). The intended
regime is small \(d\), very large \(N\), and a Fourier grid that fits in memory.

A fast matrix-vector product does not by itself give a fast solve. When the
regularization parameter \(\lambda\) is small, the eigenvalues of \(A\) may
span many orders of magnitude, and CG may take thousands of iterations. Each
iteration remains inexpensive, but their accumulated cost can dominate both
Fourier setup and prediction. The computational question studied here is
therefore narrow: how can the Fourier system be preconditioned without
discarding its fast Toeplitz product?

The diagonal structure of \(D\) gives a useful starting point. Fourier index
\(j\) contributes \(N|D_{jj}|^2\) to the diagonal of the unregularized system.
The ratio \(\rho_j=N|D_{jj}|^2/\lambda\) compares that contribution with the
ridge term. A large ratio means that the unregularized diagonal contribution is
not dominated by regularization, so the coordinate is worth treating
accurately. This is only a selection rule:
off-diagonal entries of \(G\) still couple the coordinates, so the ratio by
itself does not prove that the remaining system is well conditioned.

We select the largest-ratio Fourier indices and enclose them in a Cartesian
box \(B\). The box has two practical roles. First, it gives a contiguous
subgrid whose matrix \(A_{BB}\) can be obtained from the already computed
Toeplitz data. Second, it leaves the complement available for inexpensive
diagonal scaling. On \(B\), we use either the inverse of \(A_{BB}\) or a
rank-\(q\) correction built from its leading eigenpairs. The latter is the
spectral preconditioner used by EigenPro; here it is used inside PCG, without
EigenPro's Richardson or stochastic training iteration. We denote the two
choices by \(P_{\mathrm{inv}}(B)\) and \(P_{\mathrm{eig}}(B,q)\). The inverse
gives the stronger correction when \(|B|\) is moderate, whereas
\(P_{\mathrm{eig}}(B,q)\) can treat a larger box at limited rank. Taking
\(B=J_m\) gives the corresponding full-grid cases, not separate methods.

The analysis makes explicit what this selection rule can and cannot ensure.
The condition-number bound separates three effects: the accuracy of the
operation on \(A_{BB}\), the error of diagonal scaling on the complement, and
the coupling between the two blocks. The ratio \(\rho_j\) helps control which
coordinates remain outside the box, but the last two effects still depend on
the Toeplitz matrix generated by the sample locations. This separation links
the proof directly to the algorithm and explains why a larger box may reduce
the iteration count while increasing total time.

\subsection{Related work}
\label{subsec:related-work}

Bochner's theorem gives a spectral representation for a stationary
positive-definite kernel
\cite{bochner1933monotone,wendland2004scattered}. Random Fourier features draw
frequencies from this density and obtain a finite feature model
\cite{rahimi2007random}. The construction used here instead evaluates the
spectrum on an equispaced grid. This deterministic grid supports a controlled
Fourier approximation and makes \(G\) Toeplitz
\cite{greengard2025efgp,barnett2024uniform}. The present work leaves this
Fourier approximation fixed and addresses the conditioning of its weight-space
solve.

At the linear-solver level, diagonal (Jacobi) scaling is the simplest
preconditioner for a positive-definite system \cite{saad2003iterative}.
For positive-definite Toeplitz systems, circulant preconditioners have
long been combined with conjugate gradients and FFT-based matrix-vector
products \cite{chan1989toeplitz}. Chan also proposed an optimal
circulant approximation in the Frobenius norm
\cite{chan1988optimal}. These results do not apply directly to our
system: although \(G\) is Toeplitz, the Fourier KRR matrix
\[
    A = DGD + \lambda I
\]
is generally not Toeplitz because \(D\) is nonconstant. Our
preconditioner instead preserves the fast Toeplitz product with \(G\)
while using the Fourier weights in \(D\) to identify the coordinates
that receive a stronger correction.


Another line of work reduces the model dimension through Nystr\"om
columns or inducing locations \cite{williams2001nystrom,titsias2009variational}.
KISS-GP places a regular interpolation grid in input space
\cite{wilson2015kissgp}; the grid here instead lies in frequency space.
FALKON combines a sampled model with a preconditioned iterative solve
\cite{rudi2017falkon,meanti2020kernelroof}. Randomized Nystr\"om, sketching,
and Cholesky constructions have also been used to precondition kernel
systems \cite{avron2017faster,frangella2023randomized,
chen2022rpcholesky,diaz2023robust}. In particular, Nystr\"om
approximations can reduce the cost of spectral preconditioning while
retaining much of its effect \cite{abedsoltan2024nystrom}.

These methods typically construct data-space or inducing-point
preconditioners from sampled kernel columns or sketches.  In the end-to-end
comparison, Nystr\"om KRR and RPCholesky KRR are therefore treated as complete
data-space KRR pipelines, including their own sampling, construction, and
solve costs.  Our active-box construction instead starts after the
equispaced-Fourier model has been fixed. The resulting system has dimension
\(M\), and its Toeplitz generator is already available after Fourier setup.
The active-box preconditioner therefore uses information already present in
the Fourier system and requires no additional pass over the \(N\) training
samples.  Fourier-system adaptations of Nystr\"om or RPCholesky are not used as
the main end-to-end KRR baselines.

The location of the selected information changes the available computation.
A data-space column corresponds to a sample location and requires kernel
values against the data. Here a Cartesian frequency box gives a principal
block that can be assembled from the Toeplitz generator after Fourier setup,
without another pass over the \(N\) samples. Its subsequent costs depend on
\(|B|\), not directly on \(N\). The price is that the score uses only
\(D\), \(N\), and \(\lambda\); it cannot see the off-diagonal entries of
the sample-dependent matrix \(G\). The box is therefore a cheap proposal,
not a substitute for data-dependent spectral information. The bound retains
the missing coupling terms, and the experiments test a small set of boxes.

EigenPro introduced a spectral preconditioner that reduces the largest
eigenvalues through a correction based on leading eigenvectors
\cite{ma2017eigenpro,ma2019eigenpro,abedsoltan2023toward}.
Recent work has reduced the cost of this spectral information using
Nystr\"om approximations \cite{abedsoltan2024nystrom}, while EigenPro4
reduces the projection overhead of large kernel models through delayed
projections \cite{abedsoltan2025delayed}. These methods address spectral
preconditioning and projection in general kernel training. Our
\(P_{\mathrm{eig}}(B,q)\) restricts the spectral correction to a fixed Fourier
submatrix. This changes its storage and application cost from \(O(Mq)\) to
\(O(|B|q)\); when \(B=J_m\), it gives the corresponding full-grid correction.


Standard preconditioning theory explains CG through the spectrum of the
preconditioned matrix \cite{hestenes1952methods,saad2003iterative}. A small
condition number is not sufficient for lower wall time: the setup and repeated
application of the preconditioner must also be cheaper than the iterations they
save.  Stage~1 therefore charges every KRR pipeline for its own setup and solve.
Stage~2 holds \(A\) and \(\bvec\) identical, reports the shared Fourier setup
separately, and compares solver totals consisting of score selection (when
applicable), preconditioner construction, and CG/PCG.

\subsection{Contributions}
\label{subsec:main-contributions}

The main contributions are as follows.
\begin{enumerate}
    \item We distinguish a regularization-aware score threshold from a
    top-\(s\) selection budget. Exact threshold radii show polylogarithmic
    active-box growth for the squared exponential spectrum and algebraic
    growth for the Mat\'ern spectrum before the finite grid saturates.

    \item We combine \(P_{\mathrm{inv}}(B)\) or
    \(P_{\mathrm{eig}}(B,q)\) with diagonal scaling outside the box. A
    spectral bound separates the box correction, the diagonal approximation,
    and the interaction between them. A population Fourier-Gram formula
    further separates input-distribution correlations from finite-sample error
    in the last two terms.

    \item We specify a two-stage evaluation that does not mix model setup with
    solver effects.  Stage~1 first compares EFGP-CG with explicit inverse and
    active-box EigenPro branches---including the full-grid eigenpair limit---
    over \(10^7\) to \(3\times10^8\) samples.  Box/rank, data-set, and kernel
    changes use the same two-family reporting rule.  Data-space Nystr\"om,
    data-space RPCholesky, and EFGP+Jacobi then provide a secondary complete-KRR
    reference, with method-specific setup and solve accounting and explicit
    time--accuracy reporting.  Stage~2
    fixes \(A\) and \(\bvec\), uses paired repeated timing, and compares CG,
    Jacobi and the proposed full-grid, active-inverse, and active-eigenpair
    family members by
    total time including score selection, preconditioner construction, and
    solve.
\end{enumerate}

\Cref{sec:background} derives the Fourier system and its fast product.
\Cref{sec:main-algorithms} gives the two preconditioners, their spectral bound,
and their costs. \Cref{sec:experiments} reports the timing and accuracy results,
and \cref{sec:discussion} states the practical limits. Proofs and full
algorithms are in the appendices.

\subsection{Notation}
\label{subsec:notation}

The symbol \((\cdot)^*\) denotes conjugate transpose, and \(\|\cdot\|_2\)
denotes the Euclidean or spectral norm. For an index set \(B\), \(A_{BB}\)
is the principal submatrix indexed by \(B\). For a Hermitian positive-definite
matrix \(A\), its condition number is
\(\kappa(A)=\lambda_{\max}(A)/\lambda_{\min}(A)\), and
\(\|v\|_A=(v^*Av)^{1/2}\).
